%=====================================================================
\section{Data}
\label{sec:data}
One network is used per type. The three networks were chosen to differ in edge semantics,
generating mechanism and temporal resolution; their properties are summarised in
\cref{tab:data}.

\begin{table}[t]
\centering
\small
\setlength{\tabcolsep}{3pt}
\begin{tabular*}{\columnwidth}{@{\extracolsep{\fill}}llrrll@{}}
\toprule
Type & Dataset & $N$ & Events & Span & Res. \\
\midrule
Contact & Primary school (PS) & 242 & 125\,773 & 2 d & 20 s \\
Comm. & CollegeMsg (CM) & 1\,899 & 59\,798 & 193 d & 1 s \\
Citation & DBLP--DB (CIT) & 19\,150 & 73\,269 & 50 yr & 1 yr \\
\bottomrule
\end{tabular*}
\caption{Datasets. $N$: nodes; Res.: temporal resolution. CIT spans publication years
1970--2020; the primary analysis covers 1990--2020. PS is undirected and has class labels;
CM is directed and unlabelled; CIT is directed and acyclic.}
\label{tab:data}
\end{table}

\paragraph{Face-to-face contacts.}
The primary-school data \citep{stehle2011high,gemmetto2014mitigation} were collected by
the SocioPatterns collaboration with wearable proximity sensors over two school days. A
contact is recorded when two individuals are face to face within approximately 1.5\,m for
at least 20\,s. Each individual carries a class label, which serves as ground truth for
community detection. The data are released under CC BY-NC-SA 3.0. Activity is confined to
school hours (approximately 08:30--17:20), with up to 2,168 contact records per 10\,min and
no records overnight.

\paragraph{Online communication.}
CollegeMsg \citep{panzarasa2009patterns} contains time-stamped private messages between
1\,899 users of an online student community over 193 days at 1\,s resolution; removing
37 duplicate triples leaves 59,798 messages. Edges are directed, and no node labels are
available. Activity is strongly nonstationary: 37,680 messages (63\%) fall in May 2004, and
monthly volume declines steadily thereafter to 816 messages in October.

\paragraph{Scientific publications.}
The citation network is an induced subgraph of AMiner DBLP-Citation V14
\citep{tang2008arnetminer}, restricted to a fixed whitelist of database venues (SIGMOD,
VLDB/PVLDB, ICDE, PODS, EDBT, ICDT, TODS, VLDB Journal, Data \& Knowledge Engineering and
Information Systems) and their declared name variants. Of 5,259,858 streamed records,
19,150 publications from 1970--2020 are retained, including 6,709 isolates. A citation
$u\to v$ is dated by the publication year of $u$, and the primary graph requires
$y_v<y_u$; it contains 73,269 distinct citations and is therefore acyclic. The 5,486
same-year references, whose order cannot be resolved, are kept separately for a
sensitivity analysis; 6,390 references to later papers, 494 self-loops and 37,440
repeated reference entries are removed. Author self-citations are retained. All counts
refer to citations within the selected venues and not to global citation impact. AMiner
permits research use; no further licence is assumed. The number of retained citations per
year declines in the most recent years although the number of publications continues to
grow, which indicates coverage and censoring effects and motivates a sensitivity check
with a 2015 cutoff.

\paragraph{Preprocessing.}
All data are converted to a common event format $(u,v,t)$ with integer timestamps
(seconds or years) and a table of node labels. Self-loops and exact duplicate events are
removed. Each source has known limitations: the citation data have annual resolution,
the contact data cover only two days, and the sensors record only participating
individuals. The datasets differ by design in edge semantics (proximity, message,
citation), directedness and time scale, so that agreement across them cannot be
attributed to a single data type.
